\chapter{FINANCIAL PLAN}
\label{ch:finance}

\section{Basis of the model}

The figures in this chapter are a \textbf{projection built on stated
assumptions}, not measurements. Where a price is public it is cited; where a
quantity is a business hypothesis it is labelled as one. The assumptions are
listed first so that a reader can disagree with a number without having to
reverse-engineer where it came from.

\begin{table}[htbp]
\centering
\caption{Assumptions underlying the model}
\label{tab:assumptions}
\small
\begin{tabularx}{\textwidth}{@{}p{4.3cm} X@{}}
\toprule
\textbf{Assumption} & \textbf{Basis} \\
\midrule
Currency and conversion
  & All figures in Algerian dinars (DZD). Infrastructure is billed in foreign
    currency and converted at 1~USD~$=$~130~DZD, the 2026
    average~\cite{xe2026}. \\
\addlinespace
No AI inference cost
  & Embeddings are computed locally on open pre-trained models running on CPU
    (Chapter~\ref{ch:prototype}). There is no per-request charge to any external
    provider --- the dominant recurring cost in comparable AI products, and
    zero here. \\
\addlinespace
No software licences
  & The stack is entirely open-source (PostgreSQL, Redis, FastAPI, Next.js). \\
\addlinespace
Year~1 is unsalaried
  & The founding team works without salary during the pilot, as is normal for a
    labelled startup in its first year. Salaried hiring begins in year~2. \\
\addlinespace
Salary basis
  & Developer gross taken at 1.8~M~DZD per year, within the reported range for
    software developers in Algeria~\cite{worldsalaries2025}; a 26\,\% employer
    social charge is applied on top. \\
\addlinespace
Revenue is institutional
  & The service is free for citizens permanently (Chapter~\ref{ch:market}).
    Revenue comes from institutional subscriptions and partnerships. \\
\bottomrule
\end{tabularx}
\end{table}

\towrite{Replace the salary basis and the subscription price with real quotes
before the defence --- an actual pro-forma from a hosting provider and a real
conversation with one institution are worth more to a jury than any published
average. Everything else in this chapter recomputes from those two numbers.}

\section{Costs and expenses}

\subsection{Infrastructure}

The architecture makes this section unusually short, and that is the single most
important financial fact about the project. Because inference is local and the
stack is open-source, running the platform costs the price of one modest server,
one managed database and some object storage.

\begin{table}[htbp]
\centering
\caption{Infrastructure costs, monthly}
\label{tab:infra-costs}
\small
\begin{tabularx}{\textwidth}{@{}X r r l@{}}
\toprule
\textbf{Line} & \textbf{Y1 (USD)} & \textbf{Y2 (USD)} & \textbf{Reference} \\
\midrule
Application compute (API + worker) & 35 & 70  & Hetzner CCX23~\cite{selfhost2026} \\
Managed PostgreSQL + \code{pgvector} & 20 & 40  & Railway~\cite{railway2026} \\
Object storage and bandwidth        & 10 & 40  & S3-compatible, per GB \\
Domain and certificates             &  2 &  2  & TLS free via Let's Encrypt \\
Transactional e-mail                &  0 & 15  & free tier at pilot volume \\
Monitoring and logs                 &  0 & 10  & free tier at pilot volume \\
\midrule
\textbf{Total per month}            & \textbf{67}  & \textbf{177} & \\
\textbf{Total per year (USD)}       & \textbf{804} & \textbf{2\,124} & \\
\textbf{Total per year (DZD)}       & \textbf{104\,520} & \textbf{276\,120} & at 130 DZD/USD \\
\bottomrule
\end{tabularx}
\end{table}

Two implementation decisions hold these figures down. Uploaded photographs are
re-encoded to WebP at a bounded resolution, cutting a typical phone image by
roughly 95\,\% and postponing the point at which storage dominates. And the API
runs without Redis when the queue is disabled, so the pilot can be deployed as a
single service before the worker is needed.

\subsection{Personnel}

Personnel is the dominant cost from year~2 onward, and it follows the headcount
of Table~\ref{tab:headcount} in Chapter~\ref{ch:organization}.

\begin{table}[htbp]
\centering
\caption{Personnel expenses, in million DZD per year (gross, including 26\,\%
employer social charges)}
\label{tab:personnel-costs}
\small
\begin{tabular}{@{}l r r r r r@{}}
\toprule
\textbf{Role} & \textbf{Y1} & \textbf{Y2} & \textbf{Y3} & \textbf{Y4} & \textbf{Y5} \\
\midrule
Full-stack developer        & --- & 4.54 & 6.80 & 9.07 & 11.34 \\
Machine-learning engineer   & --- & 2.77 & 2.77 & 2.77 & 5.54 \\
Moderation and support      & --- & 1.26 & 2.52 & 3.78 & 5.04 \\
Partnerships and business   & --- & 1.76 & 1.76 & 3.53 & 3.53 \\
Marketing and communication & --- & ---  & 1.51 & 1.51 & 3.02 \\
\midrule
\textbf{Total}              & \textbf{0.00} & \textbf{10.33} & \textbf{15.36} & \textbf{20.66} & \textbf{28.47} \\
\bottomrule
\end{tabular}
\end{table}

\subsection{Consolidated costs}

\begin{table}[htbp]
\centering
\caption{Summary of charges and costs, in million DZD per year}
\label{tab:costs}
\small
\begin{tabularx}{\textwidth}{@{}X r r r r r@{}}
\toprule
\textbf{Cost line} & \textbf{Y1} & \textbf{Y2} & \textbf{Y3} & \textbf{Y4} & \textbf{Y5} \\
\midrule
Infrastructure                       & 0.10 & 0.28 & 0.55 & 0.90 & 1.50 \\
Personnel (incl.\ social charges)    & 0.00 & 10.33 & 15.36 & 20.66 & 28.47 \\
Marketing and communication          & 0.15 & 0.50 & 1.00 & 1.50 & 2.00 \\
Legal, administrative, accounting    & 0.10 & 0.20 & 0.30 & 0.40 & 0.50 \\
Tangible fixed assets (equipment)    & 0.40 & 0.20 & 0.30 & 0.40 & 0.50 \\
\midrule
\textbf{Total annual cost}           & \textbf{0.75} & \textbf{11.51} & \textbf{17.51} & \textbf{23.86} & \textbf{32.97} \\
\bottomrule
\end{tabularx}
\end{table}

The shape of this table is worth stating explicitly: \textbf{the technology is
not the expense --- the team is}. Infrastructure never exceeds 5\,\% of the
annual cost in any year of the model. A jury may reasonably ask why an
AI-powered platform is so cheap to run, and the answer is architectural, not
optimistic: an external embedding API billed per request would have made the
free-to-citizen model impossible at any meaningful volume.

\section{Revenue}

\subsection{Revenue model}

The service is free for citizens on both sides, permanently, and that is a
constraint rather than a promotion. Charging the person who has just lost their
identity papers would suppress the loss reports the engine needs; charging the
finder would suppress the scarcer side of the market outright. Revenue therefore
comes from the actors who derive institutional value from the same pool:

\begin{itemize}[leftmargin=1.4em,itemsep=4pt]
  \item \textbf{Institutional subscriptions} --- an annual account for a
        university, transport operator, airport, municipality or commercial
        centre, replacing a paper register with a searchable one and giving
        automatic matching against national loss reports. Modelled at
        60\,000~DZD per institution per year in years~1--3, rising to
        80\,000~DZD as the pool becomes more valuable.
  \item \textbf{Sponsorship and partnerships} --- from year~3, with insurers,
        telecom operators and banks, for whom association with a recovery
        service is directly relevant to their own customers.
  \item \textbf{Anonymised recovery analytics} --- aggregate statistics on what
        is lost, where and when, useful to transport operators and municipal
        services for planning. No personal data leaves the platform.
\end{itemize}

\towrite{Decide whether you also want optional paid visibility for individual
reports. It is the obvious fourth line, and it is deliberately excluded here:
letting a citizen pay to rank higher undermines the claim that suggestions are
ranked on evidence, which is the product's central promise. If you include it,
be ready to defend that tension.}

\subsection{Projections}

\begin{table}[htbp]
\centering
\caption{Revenue projection, in million DZD per year}
\label{tab:revenue}
\small
\begin{tabularx}{\textwidth}{@{}X r r r r r@{}}
\toprule
\textbf{Revenue line} & \textbf{Y1} & \textbf{Y2} & \textbf{Y3} & \textbf{Y4} & \textbf{Y5} \\
\midrule
\multicolumn{6}{@{}l}{\textit{Optimistic scenario}}\\
Institutional subscriptions      & 0.30 & 2.40 & 7.20 & 16.00 & 28.00 \\
\quad {\footnotesize partner institutions} & {\footnotesize 5} & {\footnotesize 40} & {\footnotesize 120} & {\footnotesize 200} & {\footnotesize 350} \\
Sponsorship and partnerships     & ---  & 0.60 & 3.50 & 7.50 & 13.00 \\
Analytics                        & ---  & ---  & 1.30 & 2.50 & 4.00 \\
\textbf{Total}                   & \textbf{0.30} & \textbf{3.00} & \textbf{12.00} & \textbf{26.00} & \textbf{45.00} \\
\addlinespace
\multicolumn{6}{@{}l}{\textit{Pessimistic scenario}}\\
Institutional subscriptions      & 0.10 & 1.20 & 3.50 & 8.00 & 15.00 \\
\quad {\footnotesize partner institutions} & {\footnotesize 2} & {\footnotesize 20} & {\footnotesize 58} & {\footnotesize 100} & {\footnotesize 187} \\
Sponsorship and partnerships     & ---  & ---  & 1.00 & 3.00 & 7.00 \\
\textbf{Total}                   & \textbf{0.10} & \textbf{1.20} & \textbf{4.50} & \textbf{11.00} & \textbf{22.00} \\
\bottomrule
\end{tabularx}
\end{table}

\section{Break-even analysis}

\begin{figure}[htbp]
\centering
\begin{tikzpicture}
\begin{axis}[
    width=0.94\textwidth, height=7.2cm,
    xlabel={Year}, ylabel={Million DZD},
    xtick={1,2,3,4,5},
    xticklabels={Y1,Y2,Y3,Y4,Y5},
    ymin=0, ymax=48,
    grid=major,
    grid style={draw=black!12},
    legend pos=north west,
    legend style={font=\scriptsize, draw=rulegray, fill=white},
    tick label style={font=\scriptsize},
    label style={font=\scriptsize},
  ]
  \addplot[thick, draw=todoborder, mark=*, mark size=1.7pt]
    coordinates {(1,0.75) (2,11.51) (3,17.51) (4,23.86) (5,32.97)};
  \addlegendentry{Total annual cost}

  \addplot[thick, draw=uamobblue, mark=square*, mark size=1.7pt]
    coordinates {(1,0.30) (2,3.00) (3,12.00) (4,26.00) (5,45.00)};
  \addlegendentry{Revenue --- optimistic}

  \addplot[thick, draw=uamobblue!55, dashed, mark=triangle*, mark size=2.2pt]
    coordinates {(1,0.10) (2,1.20) (3,4.50) (4,11.00) (5,22.00)};
  \addlegendentry{Revenue --- pessimistic}
\end{axis}
\end{tikzpicture}
\caption{Annual costs against the two revenue scenarios. The optimistic case
crosses into annual profitability during year~4; the pessimistic case does not
reach it within the five-year horizon.}
\label{fig:breakeven}
\end{figure}

Two conclusions follow, and the second matters more than the first.

\paragraph{The optimistic scenario reaches annual break-even in year~4}, when
revenue of 26.0~M~DZD passes costs of 23.9~M~DZD. Cumulatively the venture
remains negative until year~5, so external financing is required for roughly
four years --- which is what the ANADE mechanisms, reaching up to
10~M~DZD~\cite{anade2025}, and the support attached to the national Startup
Label~\cite{startupdz} are designed to bridge.

\paragraph{The pessimistic scenario does not break even within the horizon}, and
the gap is instructive rather than fatal. It is driven almost entirely by the
rate of institutional adoption, not by the cost of running the platform ---
infrastructure stays below 5\,\% of costs in both scenarios. The project's
financial risk is therefore a \emph{commercial} risk, concentrated in whether
partner institutions sign, and not a technical or infrastructural one. That is
where the effort of year~1 belongs, and it is why Chapter~\ref{ch:market} puts
partnerships ahead of user acquisition.

\towrite{If your jury expects a formal break-even point in units, add the
calculation: fixed costs divided by the contribution margin per institutional
subscription. With costs dominated by salaries and a subscription price of
60\,000~DZD, the year-2 break-even is roughly 190 institutions --- a figure
worth stating explicitly, because it makes the commercial challenge concrete.}
